\subsection{An explicit scalar comparison for the original arithmetic quotient volumes}\label{endpoint-volume_connection_original--an-explicit-scalar-comparison-for-the-original-arithmetic-quotient-volumes}

Date: 13 September 2026. This is the independent scalar comparison requested by the arithmetic volume upper-estimate lane. All inequalities below concern the original source and its original monic quotient. The displayed coordinate maps are invertible maps used in the proof; their Jacobians, phases, and the quotient units are retained.

\subsubsection{1. Original data and the finite answer}\label{endpoint-volume_connection_original--original-data-and-the-finite-answer}

Let the fixed packet be a complete off-line quartet ENDPOINTSOURCED7B38A9B1D61C2AFFA0A1722SLOT0END, with ENDPOINTSOURCED7B38A9B1D61C2AFFA0A1722SLOT1END and common order ENDPOINTSOURCED7B38A9B1D61C2AFFA0A1722SLOT2END. Put

ENDPOINTSOURCED7B38A9B1D61C2AFFA0A1722SLOT3END

Keep the given monic annihilator ENDPOINTSOURCED7B38A9B1D61C2AFFA0A1722SLOT4END of degree ENDPOINTSOURCED7B38A9B1D61C2AFFA0A1722SLOT5END, with all its original local lengths. Its roots in the chart ENDPOINTSOURCED7B38A9B1D61C2AFFA0A1722SLOT6END have absolute value at most ENDPOINTSOURCED7B38A9B1D61C2AFFA0A1722SLOT7END. The source is

ENDPOINTSOURCED7B38A9B1D61C2AFFA0A1722SLOT8END

The proved arithmetic density estimate supplies its actual constants ENDPOINTSOURCED7B38A9B1D61C2AFFA0A1722SLOT9END, ENDPOINTSOURCED7B38A9B1D61C2AFFA0A1722SLOT10END, ENDPOINTSOURCED7B38A9B1D61C2AFFA0A1722SLOT11END, and ENDPOINTSOURCED7B38A9B1D61C2AFFA0A1722SLOT12END:

ENDPOINTSOURCED7B38A9B1D61C2AFFA0A1722SLOT13END

Choose any ENDPOINTSOURCED7B38A9B1D61C2AFFA0A1722SLOT14END. These are the established analytic inputs of the supplied Arithmetic Endpoint Bounds note, rather than added estimates on its unknown coefficients. Define the following literal finite numbers:

ENDPOINTSOURCED7B38A9B1D61C2AFFA0A1722SLOT15END

Here ENDPOINTSOURCED7B38A9B1D61C2AFFA0A1722SLOT16END, so ENDPOINTSOURCED7B38A9B1D61C2AFFA0A1722SLOT17END, and all denominators are strictly positive. The formula for ENDPOINTSOURCED7B38A9B1D61C2AFFA0A1722SLOT18END uses the usual ENDPOINTSOURCED7B38A9B1D61C2AFFA0A1722SLOT19END convention, although the quartet has ENDPOINTSOURCED7B38A9B1D61C2AFFA0A1722SLOT20END.

For every source degree ENDPOINTSOURCED7B38A9B1D61C2AFFA0A1722SLOT21END, let ENDPOINTSOURCED7B38A9B1D61C2AFFA0A1722SLOT22END be the positive definite Gram matrix in the original ordered monomials ENDPOINTSOURCED7B38A9B1D61C2AFFA0A1722SLOT23END, let ENDPOINTSOURCED7B38A9B1D61C2AFFA0A1722SLOT24END be the coefficient matrix of remainder modulo the original ENDPOINTSOURCED7B38A9B1D61C2AFFA0A1722SLOT25END, and set

ENDPOINTSOURCED7B38A9B1D61C2AFFA0A1722SLOT26END

The finite Hermitian comparison is

ENDPOINTSOURCED7B38A9B1D61C2AFFA0A1722SLOT27END

The proof below establishes every algebraic and comparison step in (VC.6). In particular it proves ENDPOINTSOURCED7B38A9B1D61C2AFFA0A1722SLOT28END, since the matrices in (VC.6) are positive definite.

\subsubsection{2. Exact source, quotient, and raw-jet maps}\label{endpoint-volume_connection_original--exact-source-quotient-and-raw-jet-maps}

For each finite degree define the complex-linear isomorphism

ENDPOINTSOURCED7B38A9B1D61C2AFFA0A1722SLOT29END

Its inverse is ENDPOINTSOURCED7B38A9B1D61C2AFFA0A1722SLOT30END. A change of real variable gives the equality of the original norms

ENDPOINTSOURCED7B38A9B1D61C2AFFA0A1722SLOT31END

Thus the target measure in this equality has mass exactly ENDPOINTSOURCED7B38A9B1D61C2AFFA0A1722SLOT32END. There is no mass divisor in the map. Define

ENDPOINTSOURCED7B38A9B1D61C2AFFA0A1722SLOT33END

Its leading coefficient is one, its roots are the original ENDPOINTSOURCED7B38A9B1D61C2AFFA0A1722SLOT34END-chart roots divided by ENDPOINTSOURCED7B38A9B1D61C2AFFA0A1722SLOT35END, with the same multiplicities, and their absolute values are at most ENDPOINTSOURCED7B38A9B1D61C2AFFA0A1722SLOT36END. For every polynomial ENDPOINTSOURCED7B38A9B1D61C2AFFA0A1722SLOT37END,

ENDPOINTSOURCED7B38A9B1D61C2AFFA0A1722SLOT38END

The nonzero scalar ENDPOINTSOURCED7B38A9B1D61C2AFFA0A1722SLOT39END is the exact unit connecting the two displayed generators. Consequently (VC.7) induces a quotient isomorphism

ENDPOINTSOURCED7B38A9B1D61C2AFFA0A1722SLOT40END

Division is compatible with this map. Indeed, if ENDPOINTSOURCED7B38A9B1D61C2AFFA0A1722SLOT41END with ENDPOINTSOURCED7B38A9B1D61C2AFFA0A1722SLOT42END, then (VC.10) gives ENDPOINTSOURCED7B38A9B1D61C2AFFA0A1722SLOT43END, and ENDPOINTSOURCED7B38A9B1D61C2AFFA0A1722SLOT44END; uniqueness of division proves

ENDPOINTSOURCED7B38A9B1D61C2AFFA0A1722SLOT45END

If ENDPOINTSOURCED7B38A9B1D61C2AFFA0A1722SLOT46END is a root in the original ENDPOINTSOURCED7B38A9B1D61C2AFFA0A1722SLOT47END chart and ENDPOINTSOURCED7B38A9B1D61C2AFFA0A1722SLOT48END, the raw derivative of order ENDPOINTSOURCED7B38A9B1D61C2AFFA0A1722SLOT49END obeys

ENDPOINTSOURCED7B38A9B1D61C2AFFA0A1722SLOT50END

This follows by differentiating the affine composition ENDPOINTSOURCED7B38A9B1D61C2AFFA0A1722SLOT51END times. In local Taylor coordinates ENDPOINTSOURCED7B38A9B1D61C2AFFA0A1722SLOT52END, the multiplier is still ENDPOINTSOURCED7B38A9B1D61C2AFFA0A1722SLOT53END, while the factorial stays ENDPOINTSOURCED7B38A9B1D61C2AFFA0A1722SLOT54END. The derivative phase has therefore been specified for every retained local length.

For precision at the matrix level, let ENDPOINTSOURCED7B38A9B1D61C2AFFA0A1722SLOT55END be the coefficient matrix of ENDPOINTSOURCED7B38A9B1D61C2AFFA0A1722SLOT56END in increasing monomial degrees. Its columns are the coefficient vectors of ENDPOINTSOURCED7B38A9B1D61C2AFFA0A1722SLOT57END, ENDPOINTSOURCED7B38A9B1D61C2AFFA0A1722SLOT58END; its diagonal is ENDPOINTSOURCED7B38A9B1D61C2AFFA0A1722SLOT59END. Thus it is invertible and

ENDPOINTSOURCED7B38A9B1D61C2AFFA0A1722SLOT60END

Let tildes denote the target-coordinate Gram and remainder matrices. Equations (VC.8) and (VC.12) give

ENDPOINTSOURCED7B38A9B1D61C2AFFA0A1722SLOT61END

The third equality follows by inserting the first two into ENDPOINTSOURCED7B38A9B1D61C2AFFA0A1722SLOT62END. Hence ENDPOINTSOURCED7B38A9B1D61C2AFFA0A1722SLOT63END and ENDPOINTSOURCED7B38A9B1D61C2AFFA0A1722SLOT64END. The exact determinant factor is independent of ENDPOINTSOURCED7B38A9B1D61C2AFFA0A1722SLOT65END and cancels in the ratio in (VC.6). All subsequent norm estimates can therefore be pulled back to the original matrices by this proved congruence.

To state the same comparison in the original raw-jet coordinates, list the distinct roots of ENDPOINTSOURCED7B38A9B1D61C2AFFA0A1722SLOT66END as ENDPOINTSOURCED7B38A9B1D61C2AFFA0A1722SLOT67END with full lengths ENDPOINTSOURCED7B38A9B1D61C2AFFA0A1722SLOT68END, ENDPOINTSOURCED7B38A9B1D61C2AFFA0A1722SLOT69END. In the fixed ordering of pairs ENDPOINTSOURCED7B38A9B1D61C2AFFA0A1722SLOT70END, ENDPOINTSOURCED7B38A9B1D61C2AFFA0A1722SLOT71END, define the square matrix

ENDPOINTSOURCED7B38A9B1D61C2AFFA0A1722SLOT72END

Its action is the raw-derivative map on polynomials of degree below ENDPOINTSOURCED7B38A9B1D61C2AFFA0A1722SLOT73END. The exact Taylor formula shows that a polynomial has all these derivatives zero if and only if it is divisible by each ENDPOINTSOURCED7B38A9B1D61C2AFFA0A1722SLOT74END, hence by their product ENDPOINTSOURCED7B38A9B1D61C2AFFA0A1722SLOT75END. For degree below ENDPOINTSOURCED7B38A9B1D61C2AFFA0A1722SLOT76END, such a polynomial is zero. Thus ENDPOINTSOURCED7B38A9B1D61C2AFFA0A1722SLOT77END is injective, and equality of its domain and codomain dimensions proves that it is invertible. Every polynomial and its remainder differ by a multiple of ENDPOINTSOURCED7B38A9B1D61C2AFFA0A1722SLOT78END, so their listed raw derivatives agree. Therefore the degree-ENDPOINTSOURCED7B38A9B1D61C2AFFA0A1722SLOT79END raw-jet matrix is exactly ENDPOINTSOURCED7B38A9B1D61C2AFFA0A1722SLOT80END, and

ENDPOINTSOURCED7B38A9B1D61C2AFFA0A1722SLOT81END

The multiplier ENDPOINTSOURCED7B38A9B1D61C2AFFA0A1722SLOT82END is independent of source degree. Consequently (VC.6) holds with the same constant in raw-jet coordinates and has the same determinant ratio. Equation (VC.13) gives the exact phase for each component when those raw jets are expressed in the affine chart. This is the direct typed relation between the remainder-kernel estimate and the raw-jet determinant calculation.

\subsubsection{3. The Legendre coefficient bound, proved in full}\label{endpoint-volume_connection_original--the-legendre-coefficient-bound-proved-in-full}

For a polynomial ENDPOINTSOURCED7B38A9B1D61C2AFFA0A1722SLOT83END, write ENDPOINTSOURCED7B38A9B1D61C2AFFA0A1722SLOT84END. We prove

ENDPOINTSOURCED7B38A9B1D61C2AFFA0A1722SLOT85END

Define

ENDPOINTSOURCED7B38A9B1D61C2AFFA0A1722SLOT86END

The leading coefficient is ENDPOINTSOURCED7B38A9B1D61C2AFFA0A1722SLOT87END. Repeated integration by parts, with every boundary term zero because ENDPOINTSOURCED7B38A9B1D61C2AFFA0A1722SLOT88END has zeros of order ENDPOINTSOURCED7B38A9B1D61C2AFFA0A1722SLOT89END at both endpoints, proves that ENDPOINTSOURCED7B38A9B1D61C2AFFA0A1722SLOT90END is orthogonal to every polynomial of degree below ENDPOINTSOURCED7B38A9B1D61C2AFFA0A1722SLOT91END. The same integration by parts gives

ENDPOINTSOURCED7B38A9B1D61C2AFFA0A1722SLOT92END

Here the intermediate integral is ENDPOINTSOURCED7B38A9B1D61C2AFFA0A1722SLOT93END. To verify that expression without any special-function identity, substitute ENDPOINTSOURCED7B38A9B1D61C2AFFA0A1722SLOT94END, obtaining ENDPOINTSOURCED7B38A9B1D61C2AFFA0A1722SLOT95END. Repeated integration by parts proves ENDPOINTSOURCED7B38A9B1D61C2AFFA0A1722SLOT96END for nonnegative integers ENDPOINTSOURCED7B38A9B1D61C2AFFA0A1722SLOT97END: integration by parts with antiderivative ENDPOINTSOURCED7B38A9B1D61C2AFFA0A1722SLOT98END reduces ENDPOINTSOURCED7B38A9B1D61C2AFFA0A1722SLOT99END by one and increases ENDPOINTSOURCED7B38A9B1D61C2AFFA0A1722SLOT100END by one, and the case ENDPOINTSOURCED7B38A9B1D61C2AFFA0A1722SLOT101END is ENDPOINTSOURCED7B38A9B1D61C2AFFA0A1722SLOT102END. This proves (VC.18).

Leibniz's rule applied to ENDPOINTSOURCED7B38A9B1D61C2AFFA0A1722SLOT103END proves the exact finite formula

ENDPOINTSOURCED7B38A9B1D61C2AFFA0A1722SLOT104END

The coefficient norm of a product is at most the product of its coefficient norms, by the triangle inequality for each convolution coefficient. The binomial theorem gives ENDPOINTSOURCED7B38A9B1D61C2AFFA0A1722SLOT105END. Therefore

ENDPOINTSOURCED7B38A9B1D61C2AFFA0A1722SLOT106END

The middle equality is the coefficient of ENDPOINTSOURCED7B38A9B1D61C2AFFA0A1722SLOT107END in ENDPOINTSOURCED7B38A9B1D61C2AFFA0A1722SLOT108END; the last inequality bounds one coefficient by the sum of all positive coefficients. The degrees and nonzero leading coefficients show that ENDPOINTSOURCED7B38A9B1D61C2AFFA0A1722SLOT109END are a basis. Write ENDPOINTSOURCED7B38A9B1D61C2AFFA0A1722SLOT110END. By orthogonality and Cauchy--Schwarz,

ENDPOINTSOURCED7B38A9B1D61C2AFFA0A1722SLOT111END

This formula is valid for complex coefficients because ENDPOINTSOURCED7B38A9B1D61C2AFFA0A1722SLOT112END is real. Combining (VC.20)--(VC.21), bounding ENDPOINTSOURCED7B38A9B1D61C2AFFA0A1722SLOT113END by ENDPOINTSOURCED7B38A9B1D61C2AFFA0A1722SLOT114END, and summing the geometric series ENDPOINTSOURCED7B38A9B1D61C2AFFA0A1722SLOT115END proves (VC.16).

\subsubsection{4. The complete monic-division bound}\label{endpoint-volume_connection_original--the-complete-monic-division-bound}

Write the roots of ENDPOINTSOURCED7B38A9B1D61C2AFFA0A1722SLOT116END, with repetitions, as ENDPOINTSOURCED7B38A9B1D61C2AFFA0A1722SLOT117END, where ENDPOINTSOURCED7B38A9B1D61C2AFFA0A1722SLOT118END. Define the elementary symmetric polynomials ENDPOINTSOURCED7B38A9B1D61C2AFFA0A1722SLOT119END and complete homogeneous polynomials ENDPOINTSOURCED7B38A9B1D61C2AFFA0A1722SLOT120END by the finite coefficient identities

ENDPOINTSOURCED7B38A9B1D61C2AFFA0A1722SLOT121END

The second identity is a formal power-series identity: ENDPOINTSOURCED7B38A9B1D61C2AFFA0A1722SLOT122END is the sum of all monomials ENDPOINTSOURCED7B38A9B1D61C2AFFA0A1722SLOT123END with nonnegative ENDPOINTSOURCED7B38A9B1D61C2AFFA0A1722SLOT124END and ENDPOINTSOURCED7B38A9B1D61C2AFFA0A1722SLOT125END. There are ENDPOINTSOURCED7B38A9B1D61C2AFFA0A1722SLOT126END such monomials: placing ENDPOINTSOURCED7B38A9B1D61C2AFFA0A1722SLOT127END separators among ENDPOINTSOURCED7B38A9B1D61C2AFFA0A1722SLOT128END positions gives a bijection with the exponent lists. Consequently

ENDPOINTSOURCED7B38A9B1D61C2AFFA0A1722SLOT129END

For ENDPOINTSOURCED7B38A9B1D61C2AFFA0A1722SLOT130END, set ENDPOINTSOURCED7B38A9B1D61C2AFFA0A1722SLOT131END. Multiplying the two series in (VC.22) shows that every coefficient of degrees ENDPOINTSOURCED7B38A9B1D61C2AFFA0A1722SLOT132END in ENDPOINTSOURCED7B38A9B1D61C2AFFA0A1722SLOT133END is zero, while its coefficient of degree ENDPOINTSOURCED7B38A9B1D61C2AFFA0A1722SLOT134END is one. It follows that

ENDPOINTSOURCED7B38A9B1D61C2AFFA0A1722SLOT135END

The leading coefficient of ENDPOINTSOURCED7B38A9B1D61C2AFFA0A1722SLOT136END is exactly one, so removing its leading term reduces its coefficient norm by exactly one. All remaining terms are the negative of the remainder. Thus

ENDPOINTSOURCED7B38A9B1D61C2AFFA0A1722SLOT137END

For the last inequality, the convergent geometric-series product at real ENDPOINTSOURCED7B38A9B1D61C2AFFA0A1722SLOT138END gives ENDPOINTSOURCED7B38A9B1D61C2AFFA0A1722SLOT139END. For ENDPOINTSOURCED7B38A9B1D61C2AFFA0A1722SLOT140END, the remainder of ENDPOINTSOURCED7B38A9B1D61C2AFFA0A1722SLOT141END is ENDPOINTSOURCED7B38A9B1D61C2AFFA0A1722SLOT142END and its coefficient norm is one, also at most ENDPOINTSOURCED7B38A9B1D61C2AFFA0A1722SLOT143END. Linearity and the triangle inequality now give the exact finite operator estimate

ENDPOINTSOURCED7B38A9B1D61C2AFFA0A1722SLOT144END

Multiplicity was retained throughout the root list; no distinct-root interpolation or inverse Vandermonde estimate occurs here. If desired, the smaller literal finite constant

ENDPOINTSOURCED7B38A9B1D61C2AFFA0A1722SLOT145END

can replace ENDPOINTSOURCED7B38A9B1D61C2AFFA0A1722SLOT146END in (VC.4), by the preceding proof.

\subsubsection{5. Original moment control and the remainder's source norm}\label{endpoint-volume_connection_original--original-moment-control-and-the-remainders-source-norm}

For every nonnegative integer ENDPOINTSOURCED7B38A9B1D61C2AFFA0A1722SLOT147END and real ENDPOINTSOURCED7B38A9B1D61C2AFFA0A1722SLOT148END, the nonnegative Taylor terms of ENDPOINTSOURCED7B38A9B1D61C2AFFA0A1722SLOT149END imply

ENDPOINTSOURCED7B38A9B1D61C2AFFA0A1722SLOT150END

Tonelli's theorem applied to the nonnegative product density and the identity ENDPOINTSOURCED7B38A9B1D61C2AFFA0A1722SLOT151END gives

ENDPOINTSOURCED7B38A9B1D61C2AFFA0A1722SLOT152END

Consequently the actual ENDPOINTSOURCED7B38A9B1D61C2AFFA0A1722SLOT153END-th moment is at most the right side of (VC.28) integrated using (VC.29). If ENDPOINTSOURCED7B38A9B1D61C2AFFA0A1722SLOT154END has degree at most ENDPOINTSOURCED7B38A9B1D61C2AFFA0A1722SLOT155END, each monomial has absolute value at most ENDPOINTSOURCED7B38A9B1D61C2AFFA0A1722SLOT156END. This gives

ENDPOINTSOURCED7B38A9B1D61C2AFFA0A1722SLOT157END

The second inequality uses ENDPOINTSOURCED7B38A9B1D61C2AFFA0A1722SLOT158END for ENDPOINTSOURCED7B38A9B1D61C2AFFA0A1722SLOT159END, and the proved moment bound with ENDPOINTSOURCED7B38A9B1D61C2AFFA0A1722SLOT160END. On the other hand, (VC.3) on ENDPOINTSOURCED7B38A9B1D61C2AFFA0A1722SLOT161END gives

ENDPOINTSOURCED7B38A9B1D61C2AFFA0A1722SLOT162END

Set ENDPOINTSOURCED7B38A9B1D61C2AFFA0A1722SLOT163END and ENDPOINTSOURCED7B38A9B1D61C2AFFA0A1722SLOT164END. Equations (VC.8), (VC.12), (VC.16), (VC.26), (VC.30), and (VC.31), in that order, prove

ENDPOINTSOURCED7B38A9B1D61C2AFFA0A1722SLOT165END

To specify the map's type, let ENDPOINTSOURCED7B38A9B1D61C2AFFA0A1722SLOT166END be ENDPOINTSOURCED7B38A9B1D61C2AFFA0A1722SLOT167END, and let ENDPOINTSOURCED7B38A9B1D61C2AFFA0A1722SLOT168END send a class to the coefficient vector of its unique degree-below-ENDPOINTSOURCED7B38A9B1D61C2AFFA0A1722SLOT169END representative. Write ENDPOINTSOURCED7B38A9B1D61C2AFFA0A1722SLOT170END, and ENDPOINTSOURCED7B38A9B1D61C2AFFA0A1722SLOT171END, ENDPOINTSOURCED7B38A9B1D61C2AFFA0A1722SLOT172END. Then ENDPOINTSOURCED7B38A9B1D61C2AFFA0A1722SLOT173END, and the original abstract quotient map ENDPOINTSOURCED7B38A9B1D61C2AFFA0A1722SLOT174END satisfies ENDPOINTSOURCED7B38A9B1D61C2AFFA0A1722SLOT175END. The operator estimated in (VC.32) is exactly ENDPOINTSOURCED7B38A9B1D61C2AFFA0A1722SLOT176END. Its restriction to ENDPOINTSOURCED7B38A9B1D61C2AFFA0A1722SLOT177END is the identity. In the formulas that follow, coefficient vectors and polynomials are identified only by the displayed ENDPOINTSOURCED7B38A9B1D61C2AFFA0A1722SLOT178END.

\subsubsection{6. From the bounded remainder to the kernel inequality}\label{endpoint-volume_connection_original--from-the-bounded-remainder-to-the-kernel-inequality}

The positive density in (VC.3) makes every polynomial Gram matrix positive definite. The remainder map ENDPOINTSOURCED7B38A9B1D61C2AFFA0A1722SLOT179END is surjective because its restriction to degrees below ENDPOINTSOURCED7B38A9B1D61C2AFFA0A1722SLOT180END is the identity. Thus ENDPOINTSOURCED7B38A9B1D61C2AFFA0A1722SLOT181END is positive definite. For completeness, define

ENDPOINTSOURCED7B38A9B1D61C2AFFA0A1722SLOT182END

Then ENDPOINTSOURCED7B38A9B1D61C2AFFA0A1722SLOT183END by the definition of ENDPOINTSOURCED7B38A9B1D61C2AFFA0A1722SLOT184END. For ENDPOINTSOURCED7B38A9B1D61C2AFFA0A1722SLOT185END, ENDPOINTSOURCED7B38A9B1D61C2AFFA0A1722SLOT186END. Every lift of a quotient vector ENDPOINTSOURCED7B38A9B1D61C2AFFA0A1722SLOT187END is uniquely ENDPOINTSOURCED7B38A9B1D61C2AFFA0A1722SLOT188END, ENDPOINTSOURCED7B38A9B1D61C2AFFA0A1722SLOT189END, so

ENDPOINTSOURCED7B38A9B1D61C2AFFA0A1722SLOT190END

The minimum lift norm is therefore exactly ENDPOINTSOURCED7B38A9B1D61C2AFFA0A1722SLOT191END. Degree inclusion enlarges the set of lifts and gives ENDPOINTSOURCED7B38A9B1D61C2AFFA0A1722SLOT192END. Taking ENDPOINTSOURCED7B38A9B1D61C2AFFA0A1722SLOT193END in (VC.32), the degree-ENDPOINTSOURCED7B38A9B1D61C2AFFA0A1722SLOT194END remainder is a lift of the same ENDPOINTSOURCED7B38A9B1D61C2AFFA0A1722SLOT195END, so

ENDPOINTSOURCED7B38A9B1D61C2AFFA0A1722SLOT196END

Together with ENDPOINTSOURCED7B38A9B1D61C2AFFA0A1722SLOT197END, this gives ENDPOINTSOURCED7B38A9B1D61C2AFFA0A1722SLOT198END. Congruence by ENDPOINTSOURCED7B38A9B1D61C2AFFA0A1722SLOT199END bounds every eigenvalue of the middle matrix between ENDPOINTSOURCED7B38A9B1D61C2AFFA0A1722SLOT200END and one. Inversion reverses these positive eigenvalue inequalities; conjugating back proves the first inequality chain in (VC.6). Taking products of the same ENDPOINTSOURCED7B38A9B1D61C2AFFA0A1722SLOT201END eigenvalues proves its determinant statement. This proves the finite answer.

\subsubsection{7. The complete growth calculation and its exact gap}\label{endpoint-volume_connection_original--the-complete-growth-calculation-and-its-exact-gap}

Let

ENDPOINTSOURCED7B38A9B1D61C2AFFA0A1722SLOT202END

These are bounding constants for the unchanged source. Since ENDPOINTSOURCED7B38A9B1D61C2AFFA0A1722SLOT203END, ENDPOINTSOURCED7B38A9B1D61C2AFFA0A1722SLOT204END. Also ENDPOINTSOURCED7B38A9B1D61C2AFFA0A1722SLOT205END. It follows that

ENDPOINTSOURCED7B38A9B1D61C2AFFA0A1722SLOT206END

The constants in (VC.4) further satisfy

ENDPOINTSOURCED7B38A9B1D61C2AFFA0A1722SLOT207END

On ENDPOINTSOURCED7B38A9B1D61C2AFFA0A1722SLOT208END, the derivative of ENDPOINTSOURCED7B38A9B1D61C2AFFA0A1722SLOT209END is ENDPOINTSOURCED7B38A9B1D61C2AFFA0A1722SLOT210END. Integration from zero to ENDPOINTSOURCED7B38A9B1D61C2AFFA0A1722SLOT211END gives

ENDPOINTSOURCED7B38A9B1D61C2AFFA0A1722SLOT212END

Finally ENDPOINTSOURCED7B38A9B1D61C2AFFA0A1722SLOT213END, and ENDPOINTSOURCED7B38A9B1D61C2AFFA0A1722SLOT214END. Substituting (VC.37)--(VC.39) into the literal constant in (VC.4), and using ENDPOINTSOURCED7B38A9B1D61C2AFFA0A1722SLOT215END, yields

ENDPOINTSOURCED7B38A9B1D61C2AFFA0A1722SLOT216END

Thus an explicit upper bound on the requested original central volume is

ENDPOINTSOURCED7B38A9B1D61C2AFFA0A1722SLOT217END

All constants in (VC.41) depend only on the fixed packet and the chosen ENDPOINTSOURCED7B38A9B1D61C2AFFA0A1722SLOT218END; no constant depends on ENDPOINTSOURCED7B38A9B1D61C2AFFA0A1722SLOT219END or ENDPOINTSOURCED7B38A9B1D61C2AFFA0A1722SLOT220END. Since ENDPOINTSOURCED7B38A9B1D61C2AFFA0A1722SLOT221END, one has ENDPOINTSOURCED7B38A9B1D61C2AFFA0A1722SLOT222END and ENDPOINTSOURCED7B38A9B1D61C2AFFA0A1722SLOT223END. In particular

ENDPOINTSOURCED7B38A9B1D61C2AFFA0A1722SLOT224END

The proved central-window lower estimate in the companion factor-four note is

ENDPOINTSOURCED7B38A9B1D61C2AFFA0A1722SLOT225END

The precise comparison between (VC.41) and (VC.43) is therefore an interval of admissible growth, not an RH contradiction. After division by ENDPOINTSOURCED7B38A9B1D61C2AFFA0A1722SLOT226END, the leading term of the upper bound is ENDPOINTSOURCED7B38A9B1D61C2AFFA0A1722SLOT227END. This diverges like ENDPOINTSOURCED7B38A9B1D61C2AFFA0A1722SLOT228END when ENDPOINTSOURCED7B38A9B1D61C2AFFA0A1722SLOT229END, and like ENDPOINTSOURCED7B38A9B1D61C2AFFA0A1722SLOT230END when ENDPOINTSOURCED7B38A9B1D61C2AFFA0A1722SLOT231END. The upper bound is larger than the lower scale by those explicit factors. This calculation locates exactly the loss in this scalar route: (VC.16) controls every coefficient of a degree-ENDPOINTSOURCED7B38A9B1D61C2AFFA0A1722SLOT232END polynomial, and (VC.31) replaces the full source by its least value on a length-ENDPOINTSOURCED7B38A9B1D61C2AFFA0A1722SLOT233END interval. The remainder step itself has only ENDPOINTSOURCED7B38A9B1D61C2AFFA0A1722SLOT234END logarithmic cost by (VC.39).

The sharper raw-jet matrix determinant comparison in the parent lane retains the directional dependence discarded by the scalar estimates (VC.16) and (VC.31). Equations (VC.7)--(VC.15) and (VC.33)--(VC.35) give the exact maps through which the present scalar estimate bounds that same original quotient. No claim about the sharper determinant's asymptotic size is made here.

\subsubsection{8. Provenance and verification boundary}\label{endpoint-volume_connection_original--provenance-and-verification-boundary}

The inherited analytic input is (VC.3) and the finiteness of ENDPOINTSOURCED7B38A9B1D61C2AFFA0A1722SLOT235END, proved in the complete supplied \texttt{Tau\_Arithmetic\_Endpoint\_Bounds/NOTE.tex}. The central lower bound (VC.43) is proved in the owner's complete \texttt{FOUR\_VOLUME\_THRESHOLD.md} continuation. This note proves (VC.4)--(VC.42) directly from those established input theorems, with full finite calculations. It does not replace the arithmetic weight, discard the mass, erase a raw-jet phase, shorten a local length, postulate the desired upper estimate, or assert RH. Any finite symbolic test accompanying this note tests the algebraic maps and constants; the written argument proves the estimates for all the indicated degrees.
